This project investigates the motion of charged particles in prescribed
electromagnetic fields through analytic theory and numerical simulation.
The equations of motion arising from the Lorentz force are formulated as
a first-order dynamical system and integrated numerically using an
adaptive Runge--Kutta scheme (RK45) in Python via
\texttt{scipy.integrate.solve\_ivp}. The solver is verified against known
analytic results for circular gyromotion, helical motion, the
$\mathbf{E} \times \mathbf{B}$ drift, and gradient-$B$ and curvature
drifts.

The validated solver is then applied to particle dynamics in a magnetic
dipole field. Bounce motion between mirror points is demonstrated,
the magnetic moment adiabatic invariant is verified to be approximately
conserved, and a guiding-centre integrator is compared with the full
Lorentz orbit. The framework is also applied in physical SI units,
simulating a 10\,keV electron at $L = 3$ in Earth's dipole field, and
extended to tilted dipole fields representing Neptune ($47^\circ$) and
Uranus ($59^\circ$). A corotation electric field and a rotating tilted
dipole complete the progression toward a realistic planetary
magnetosphere model. Results throughout are consistent with adiabatic
particle orbit theory.
